\section{Voronoi (wersja 1.0.54)}
\label{sec:voronoi}

\subsection*{Okno i wolumen}
Okno eksperymentu: \textbf{2025-08-24 14{:}31{:}51–2025-08-24 14{:}37{:}25} ($\approx$ \SI{5.6}{\minute}). 
Zarejestrowano \textbf{465} przepływy, wolumen $\approx$ \textbf{9.61 MiB} 
(\textbf{RX $\approx$ 9.21 MiB}, \textbf{TX $\approx$ 0.39 MiB}).

\paragraph{Przebieg badania:} 
Scenariusz obejmował: \emph{nasłuch tła (2 min), utworzenie konta, logowanie, nawigację po ekranach (online/offline), wyszukiwanie, wylogowanie}.

\subsection{Wolumen i liczność wg protokołu}
\begin{table}[H]
  \centering
  \caption{Podsumowanie wg protokołu (Voronoi, ta sesja)}
  \label{tab:voronoi-proto}
  \begin{tabular}{lrrrr}
    \toprule
    \textbf{Protokół} & \textbf{Flows} & \textbf{MiB (total)} & \textbf{B wysłane} & \textbf{B odebrane} \\
    \midrule
    HTTPS & 174 & 9.53 & 373{,}137 & 9{,}622{,}074 \\
    TLS   & 226 & 0.06 & 33{,}132  & 33{,}984 \\
    DNS   & 42  & 0.01 & 1{,}881   & 3{,}620 \\
    TCP   & 23  & 0.00 & 2{,}420   & 2{,}144 \\
    \midrule
    \textbf{Razem} & \textbf{465} & \textbf{9.61} & \textbf{410{,}570} & \textbf{9{,}661{,}822} \\
    \bottomrule
  \end{tabular}
\end{table}

\paragraph{Obserwacja:} 
Ruch aplikacji jest zdominowany przez \textbf{HTTPS}; kanały \textbf{TLS}/\textbf{TCP} mają znikomy wolumen. Nie odnotowano QUIC ani jawnego HTTP.

\subsection{Największe cele (IP/port) wg wolumenu}
\begin{table}[H]
  \centering
  \caption{Największe kierunki ruchu (Voronoi, ta sesja)}
  \label{tab:voronoi-topdst}
  \begin{tabular}{lllrr}
    \toprule
    \textbf{IP docelowe} & \textbf{Port} & \textbf{Proto} & \textbf{Flows} & \textbf{MiB} \\
    \midrule
    18.245.46.92   & 443 & HTTPS & 60 & 4.37 \\
    3.160.212.30   & 443 & HTTPS & 17 & 1.40 \\
    216.58.209.3   & 443 & HTTPS & 13 & 1.22 \\
    3.98.16.223    & 443 & HTTPS & 6  & 0.33 \\
    3.160.212.123  & 443 & HTTPS & 6  & 0.32 \\
    35.182.46.15   & 443 & HTTPS & 11 & 0.28 \\
    99.79.136.123  & 443 & HTTPS & 10 & 0.18 \\
    3.97.54.246    & 443 & HTTPS & 1  & 0.15 \\
    \bottomrule
  \end{tabular}
\end{table}

\paragraph{Obserwacja:}
Największy wolumen skupia się na 443/TCP do chmury (adresy AWS/Google) — typowe dla API/CDN aplikacji.

\subsection{Najczęstsze hosty (liczba przepływów)}
\begin{table}[H]
  \centering
  \caption{Top hosty (Voronoi, liczba przepływów)}
  \label{tab:voronoi-hosts}
  \begin{tabular}{lr}
    \toprule
    \textbf{Host} & \textbf{Flows} \\
    \midrule
    \texttt{api.voronoiapp.com} & 221 \\
    \texttt{cdn.voronoiapp.com} & 182 \\
    \texttt{fonts.gstatic.com} & 15 \\
    \texttt{firebaselogging-pa.googleapis.com} & 14 \\
    \texttt{firestore.googleapis.com} & 10 \\
    \texttt{cognito-identity.ca-central-1.amazonaws.com} & 6 \\
    \texttt{crashlyticsreports-pa.googleapis.com} & 5 \\
    \texttt{cognito-idp.ca-central-1.amazonaws.com} & 4 \\
    \texttt{firebase-settings.crashlytics.com} & 2 \\
    \texttt{firebaseinstallations.googleapis.com} & 2 \\
    \bottomrule
  \end{tabular}
\end{table}

\paragraph{Obserwacja:}
Dominują domeny aplikacji (\texttt{api.voronoiapp.com}, \texttt{cdn.voronoiapp.com}); obecne są usługi Google Firebase/Crashlytics oraz Amazon Cognito (tożsamość).

\subsection{Obserwacje jakościowe}
\begin{itemize}
  \item \textbf{Brak QUIC/HTTP/3} — transport klasyczny TLS/TCP po 443.
  \item \textbf{Brak HTTP w badanym oknie czasowym} — w badanym oknie nie zaobserwowano nieszyfrowanych żądań.
  \item \textbf{Integracje zewnętrzne} — logowanie/konfiguracja przez \texttt{cognito-identity}/\texttt{cognito-idp}; telemetria do \texttt{crashlytics}; logowanie zdarzeń do \texttt{firebaselogging}.
\end{itemize}

\subsection{Wyniki z Exodus Privacy \cite{ExodusPrivacy}}
Raport dla \texttt{com.voronoi.organization.app} (wersja 1.0.34, Google Play) ujawniła obecność \textbf{3 trackerów} oraz \textbf{13 uprawnień}.

\textbf{Trackery:}
\begin{itemize}
  \item \textbf{Facebook Share SDK} — integracja społecznościowa, potencjalnie przesyłająca metadane o interakcjach,
  \item \textbf{Google CrashLytics} — monitorowanie stabilności i raportowanie awarii,
  \item \textbf{Google Firebase Analytics} — analiza użycia aplikacji i zachowań użytkowników.
\end{itemize}

\textbf{Uprawnienia — przykłady kluczowe:}
\begin{itemize}
  \item \textbf{Pamięć masowa:} \texttt{READ\_EXTERNAL\_STORAGE}, \texttt{WRITE\_EXTERNAL\_STORAGE},
  \item \textbf{Reklama/atrybucja:} \texttt{ACCESS\_ADSERVICES\_AD\_ID}, \texttt{ACCESS\_ADSERVICES\_ATTRIBUTION},
  \item \textbf{Sieć/system:} \texttt{INTERNET}, \texttt{ACCESS\_NETWORK\_STATE}, \texttt{POST\_NOTIFICATIONS}, \texttt{READ\_GSERVICES},
  \item \textbf{Inne:} \texttt{RECEIVE} (FCM push).
\end{itemize}

\emph{Spostrzeżenia:} 
\begin{itemize}
  \item Exodus bazuje na analizie statycznej APK — obecność trackerów oznacza włączenie ich bibliotek, ale nie przesądza o faktycznej aktywności w danej sesji.
  \item Zidentyfikowane trackery pełnią głównie rolę diagnostyczną i analityczną; zakres uprawnień jest stosunkowo wąski i dotyczy przede wszystkim pamięci, sieci oraz powiadomień.
\end{itemize}

\subsection{Subiektywne spostrzeżenia}
\paragraph{Obserwowalność treści.} 
Szyfrowanie TLS ogranicza wgląd do metadanych; wolumen wskazuje na standardowe korzystanie z API/CDN (brak jawnego przesyłu danych wrażliwych).

\paragraph{Uwagi dot. bezpieczeństwa.}
Zachowanie sieciowe (HTTPS, brak plaintext, integracja z Cognito) jest spójne z typową architekturą: tożsamość w AWS Cognito, treści przez CDN/API. Pełna ocena (pinning certyfikatów, obsługa błędów TLS, polityka uprawnień w runtime) wymaga analizy APK/testów dynamicznych.

\subsection{Podsumowanie}
W badanej sesji Voronoi wykazuje:
\begin{enumerate}
\item dominację HTTPS ($\approx 9.61$ MiB) w \textbf{465} flow, okno \SI{5.6}{\minute},
\item koncentrację na hostach aplikacji (\texttt{api.voronoiapp.com}, \texttt{cdn.voronoiapp.com}) i usługach Google/AWS (Firebase, Crashlytics, Cognito),
\item brak śladów QUIC i HTTP w badanym oknie czasowym.
\end{enumerate}
